%Text automatically generated by txt2latex
\section*{DESCRIPTION}
\emph{codata} is a Fortran library providing the fundamental physical
constants published by the CODATA. 
A C API allows usage from C, or can be used as an interoperability 
layer for other languages.
A Python wrapper allows easy usage from Python.
  
The constants are implemented as derived type which carries the name,
the value, the uncertainty and the unit.
All the \emph{codata} constants are provided as double precision reals.
The names are quite long and can be aliased with shorter names.

\section*{FEATURES}
It provides:
\begin{itemize}
\item \emph{codata} constants 2022
\item \emph{codata} constants 2018
\item \emph{codata} constants 2014
\item \emph{codata} constants 2010
\end{itemize}

\section*{DOCUMENTATION }
POSIX-style man-pages are generated from the sources using 
the prep(1) preprocessor and are provided with the binaries. 
The API is written as man-pages and is basically a list of the constants.
\begin{description}
\item[MAN] \href{https://milanskocic.github.io.codata/man}{milanskocic.github.io.codata/man} 
\item[PDF] \href{https://milanskocic.github.io/codata/man/codata-bookman.pdf}{milanskocic.github.io/codata/man/codata-bookman.pdf} 
\end{description}
  
The man-pages are included in the reference manual.
\begin{description}
\item[PDF] \href{https://milanskocic.github.io/codata/latex/codata-refman.pdf}{milanskocic.github.io/codata/latex/codata-refman.pdf} 
\item[HTML] \href{https://milanskocic.github.io/codata/index.html}{milanskocic.github.io/codata/index.html} 
\end{description}

\section*{NOTES FOR FORTRAN USERS}
The \emph{codata} constant from 2022 were integrated in the Fortran 
standard library (stdlib).
This library is complementary to the constants defined in the stdlib
by providing older values for the constants.
See \href{https://github.com/fortran-lang/stdlib/releases/tag/v0.7.0}{github.com/fortran-lang/stdlib/releases/tag/v0.7.0.} 
If you need older values for the CODATA constants you can use this library 
otherwise use the standard library.
  
To use \emph{codata} within your fpm(1) (\href{https://github.com/fortran-lang/fpm}{github.com/fortran-lang/fpm)} 
project, add the following lines to your file:

\begin{verbatim}
     [dependencies]
     codata = { git="https://github.com/MilanSkocic/codata.git" }

\end{verbatim}
If you only need sources for the \emph{codata} constants that you can integrate
directly embed in your sources you may be interested by this repository
\href{https://github.com/vmagnin/fundamental\_constants}{github.com/vmagnin/fundamental\_constants.} 

\section*{LICENSE}
MIT
